\chapter{Singular Values}\label{linear:ch:svd}

\begin{center}\itshape
The singular value decomposition is the Swiss Army knife of linear algebra.
\upshape(a saying among numerical analysts)
\end{center}

A photograph with a million pixels asks a practical question: how many numbers do you really need to keep? Nearby pixels often carry much the same information. A matrix can reveal that repetition, provided we let its input and output use different coordinates.

Eugenio Beltrami studied just this freedom in 1873. His object was a bilinear form $x^{\mathsf T}Ay$. Rotate the $x$-coordinates and the $y$-coordinates independently, and the form becomes $\sum_i\sigma_ix_iy_i$. The squares $\sigma_i^2$ come from $A^{\mathsf T}A$. Camille Jordan reached the same decomposition in 1874, independently. Sylvester returned to it in 1889; Schmidt's 1907 work on integral operators also gave a best-approximation theorem. Eckart and Young proved the matrix version for factor analysis in 1936. The sources at the end trace these separate discoveries.

\portrait{jordan}{Camille Jordan (1838--1922). Photograph by Eugène Pirou, between 1889 and 1922; Accademia delle Scienze di Torino. Public domain, via Wikimedia Commons.}

Computers made this old theorem a daily tool. Golub and Kahan's 1965 algorithm helped make its computation reliable. In the Netflix Prize, a large table recorded users' film ratings, with most positions empty. A low-rank model assigned a short vector to each user and each film, predicting a rating by an inner product. Simon Funk's public account helped popularise this approach during the competition. The winning team combined many models; matrix factorisation was one of its main ingredients. Predicting unseen entries needs assumptions beyond the decomposition itself, as the last exercise makes clear.

The algebra starts with a simple bargain. A square matrix may refuse to supply an eigenbasis. Allow a different orthonormal basis on each side of the map, and every rectangular matrix admits a diagonal description.

\begin{think}
You have a $1000\times1000$ grey-scale picture and room for $50{,}000$ numbers. Which information should survive? A rank-$24$ truncation stores $48{,}024$ numbers. Theorem~\ref{linear:thm:eckartyoung} proves its optimality among matrices of that rank for two precise error measures.
\end{think}

\section{The decomposition}

Take $A\in M_{m,n}(\mathbb K)$, where $\mathbb K$ is $\R$ or $\C$. Both coordinate spaces carry their usual inner products. The input has dimension $n$ and the output dimension $m$; there is no need for these dimensions to agree.

\begin{theorem}[singular value decomposition]\label{linear:thm:svd}
Suppose $A\in M_{m,n}(\mathbb K)$ has rank $r$. We can choose orthonormal bases $v_1,\ldots,v_n$ and $u_1,\ldots,u_m$ of its input and output spaces, respectively, and positive numbers $\sigma_1\ge\cdots\ge\sigma_r$ with
\[
Av_i=\sigma_iu_i\ (1\le i\le r),\qquad Av_i=0\ (r<i\le n).
\]
In matrix notation,
\[
A=U\Sigma V^*=\sum_{i=1}^r\sigma_iu_iv_i^*,
\]
where $U,V$ are unitary and the rectangular matrix $\Sigma$ has entries $\Sigma_{ii}=\sigma_i$ for $i\le r$ and zero elsewhere. The numbers $\sigma_i$, called the \emph{singular values}, are unique. Their squares are precisely the non-zero eigenvalues of each of $A^*A$ and $AA^*$, counted with multiplicity.
\end{theorem}
\begin{proof}
Start with the square matrix $A^*A$. It is Hermitian, and $x^*A^*Ax=\|Ax\|^2\ge0$. Theorem~\ref{linear:thm:spectral} therefore gives an orthonormal eigenbasis $v_1,\ldots,v_n$ with eigenvalues $s_1\ge\cdots\ge s_n\ge0$.

The identity $\|Ax\|^2=x^*A^*Ax$ proves $\ker(A^*A)=\ker A$: either kernel condition forces $Ax=0$. Rank--nullity now says that exactly $r$ eigenvalues are positive. For $i\le r$ set $\sigma_i=\sqrt{s_i}$ and $u_i=Av_i/\sigma_i$. Their inner products are
\[
\langle u_i,u_j\rangle
=\frac{\langle A^*Av_i,v_j\rangle}{\sigma_i\sigma_j}
=\frac{s_i\delta_{ij}}{\sigma_i\sigma_j}=\delta_{ij}.
\]
Complete these vectors to an orthonormal basis of the output space. For $i>r$, $\|Av_i\|^2=s_i=0$, so $AV=U\Sigma$. Multiplication by $V^*$ gives the claimed factorisation and its rank-one expansion.

Finally, the factorisation yields $A^*A=V\Sigma^*\Sigma V^*$ and $AA^*=U\Sigma\Sigma^*U^*$. Both diagonal middle factors have positive entries $\sigma_1^2,\ldots,\sigma_r^2$ and otherwise zeros. This identifies the singular values independently of all choices of vectors.
\end{proof}

Call $v_i$ a right singular vector and $u_i$ a left singular vector. Changing their paired signs, or their paired complex phases, leaves $A$ unchanged. Repeated singular values allow further orthonormal changes within the corresponding spaces. The numbers survive all these choices.

The same bases describe four useful subspaces. The first $r$ left vectors span $\operatorname{Im}A$, and the remaining left vectors span $\ker A^*$. On the input side, the first $r$ right vectors span $\operatorname{Im}A^*$, and the rest span $\ker A$. Thus the decomposition also gives the orthogonal-complement identities recorded in Proposition~\ref{linear:prop:adjoint}, directly from the bases.

\paragraph{Geometry.} The factor $V^*$ changes to input coordinates, $\Sigma$ stretches the coordinate axes, and $U$ places the stretched axes in the output space. Over $\R$ the orthogonal factors may reflect as well as rotate; over $\C$ they are unitary changes of coordinates. The unit \emph{ball} maps to the filled ellipsoid
\[
\left\{\sum_{i\le r}y_iu_i:\ \sum_{i\le r}\frac{|y_i|^2}{\sigma_i^2}\le1\right\}.
\]
Indeed, write $x=\sum c_iv_i$ and use $y_i=\sigma_ic_i$. If $\ker A=\{0\}$, the unit sphere maps to the boundary of this ellipsoid in $\operatorname{Im}A$. If the kernel is non-zero, unused kernel coordinates let the unit sphere cover the whole filled ellipsoid. The zero map gives the single point $0$.

\begin{figure}[ht]
\centering
\fig{linear-ch06-ellipse}
\caption{The matrix $A=\bigl(\begin{smallmatrix}3&0\\4&5\end{smallmatrix}\bigr)$ maps the unit circle to an ellipse with
semi-axes $\sigma_1=3\sqrt5$ and $\sigma_2=\sqrt5$ (drawn at reduced scale). The orthogonal directions $v_1,v_2$ are sent to
the orthogonal axes $\sigma_1u_1,\sigma_2u_2$; no other pair of orthogonal directions stays orthogonal.}
\label{linear:fig:ellipse}
\end{figure}

For a normal matrix, write $A=WDW^*$ and split each diagonal entry as $\lambda_i=|\lambda_i|e^{i\theta_i}$. Absorb the phases into one unitary factor. Its singular values are exactly $|\lambda_i|$. A nilpotent Jordan block shows how far this can fail without normality: its eigenvalues are all zero, while its singular values are $1,\ldots,1,0$. Exercise~\ref{linear:ex:normsvd} gives a different two-dimensional example.

We can also put a rectangular matrix inside a Hermitian one. Define its \emph{Hermitian dilation} by
\[
\mathcal H(A)=\begin{pmatrix}0&A\\A^*&0\end{pmatrix}\in M_{m+n}(\mathbb K).
\]

\begin{proposition}\label{linear:prop:dilation}
Put $p=\min(m,n)$ and append zeros to the singular-value list so that $\sigma_i=0$ for $r<i\le p$. The spectrum of $\mathcal H(A)$ consists of $\sigma_1,-\sigma_1,\ldots,\sigma_p,-\sigma_p$, together with $|m-n|$ additional zeros. In decreasing order its $k$th eigenvalue is $\sigma_k(A)$ for every $k\le p$.
\end{proposition}
\begin{proof}
The action is $\mathcal H(A)(u,v)=(Av,A^*u)$. For each $i\le p$, the vectors $(u_i,v_i)$ and $(u_i,-v_i)$ have eigenvalues $\sigma_i$ and $-\sigma_i$. This includes $i>r$, because then $Av_i=0$ and $A^*u_i=0$. Any remaining $(u_j,0)$ or $(0,v_j)$ lies in the kernel. We have produced $m+n$ non-zero, mutually orthogonal vectors, so the list is complete. Its first $p$ entries in decreasing order are the non-negative singular values, including their appended zeros.
\end{proof}

The dilation lets us use Hermitian eigenvalue inequalities on rectangular matrices. Since $\mathcal H(A+B)=\mathcal H(A)+\mathcal H(B)$, Weyl's inequality (Theorem~\ref{linear:thm:weyl}) gives
\begin{equation}\label{linear:eq:weylsv}
\sigma_{i+j-1}(A+B)\le\sigma_i(A)+\sigma_j(B),\qquad i+j-1\le p.
\end{equation}

\section{The polar decomposition}

For a complex number, modulus and phase separate as $z=e^{i\theta}|z|$. A square matrix has a corresponding separation into a positive factor and a unitary factor.

\begin{theorem}[polar decomposition]\label{linear:thm:polar}
Each $A\in M_n(\mathbb K)$ has a factorisation $A=WP$, where $W$ is unitary and $P$ is positive semidefinite. There is only one possible positive factor: $P=\sqrt{A^*A}$. If $A$ is invertible, its unitary factor is unique as well.
\end{theorem}
\begin{proof}
An SVD gives
\[
A=(UV^*)(V\Sigma V^*).
\]
The first factor is unitary, and the second is positive semidefinite. Any factorisation $A=WP$ of the stated kind satisfies $A^*A=P^2$.

Here is why a positive semidefinite square root is unique. If $P^2=S$, then $P$ commutes with $S$ and hence preserves each eigenspace of $S$. On the eigenspace for $s\ge0$, diagonalise the Hermitian restriction of $P$. Its eigenvalues are non-negative numbers whose squares equal $s$, so that restriction is $\sqrt{s}\,I$. This determines $P$ on every eigenspace of $S$. When $A$ is invertible, $P$ is invertible too, and $W=AP^{-1}$ has no remaining freedom.
\end{proof}

The positive factor stretches along orthogonal axes. The unitary factor changes their orientation. Over $\R$, if $\det A>0$, then $\det W=1$. Exercise~\ref{linear:ex:realnormal} supplies orthogonal coordinates in which $W$ consists of rotation blocks and scalar blocks $1$ and $-1$. There are an even number of $-1$ blocks; pair them as rotations by $\pi$. Turn every angle gradually from zero to its final value to get a path $W_t$ from $I$ to $W$. The matrices $P_t=(1-t)I+tP$ are positive definite when $A$ is invertible. Hence $W_tP_t$ connects $I$ to $A$ through invertible matrices of positive determinant. This proves that this set is path connected.

\section{Norms of matrices}

Three sizes of a matrix will recur. Let $p=\min(m,n)$ and include zero singular values up to index $p$.
\begin{itemize}
\item The \emph{operator norm} is $\|A\|=\max_{\|x\|=1}\|Ax\|=\sigma_1$. In singular coordinates, $\|Ax\|^2=\sum_i\sigma_i^2|c_i|^2$, which proves the formula. Applying the definition twice gives $\|AB\|\le\|A\|\|B\|$.
\item The \emph{Frobenius norm} counts all entries:
$\|A\|_F=(\sum_{i,j}|a_{ij}|^2)^{1/2}=(\operatorname{tr}A^*A)^{1/2}=(\sum_i\sigma_i^2)^{1/2}$.
\item For $1\le k\le p$, the \emph{Ky Fan $k$-norm} is $\|A\|_{(k)}=\sum_{i=1}^k\sigma_i$. At $k=p$ it is the \emph{nuclear norm}, also called the trace norm, written $\|A\|_*$.
\end{itemize}
Multiplying on either side by a unitary matrix preserves the singular values, so all these quantities are \emph{unitarily invariant}. Positivity and homogeneity of the Ky Fan quantities follow from the formulas. Their triangle inequality needs an argument.

\begin{proposition}\label{linear:prop:kyfannorm}
For $1\le k\le p$ and matrices $A,B$ of the same size,
\[
\|A+B\|_{(k)}\le\|A\|_{(k)}+\|B\|_{(k)}.
\]
\end{proposition}
\begin{proof}
Proposition~\ref{linear:prop:dilation} identifies $\|A\|_{(k)}$ with the sum of the largest $k$ eigenvalues of $\mathcal H(A)$. Exercise~\ref{linear:ex:kyfan} proves that this eigenvalue sum is subadditive on Hermitian matrices. Apply that result to the two dilations and use the linearity of $\mathcal H$.
\end{proof}

\section{The best approximation of lower rank}

Keep just the first $k$ singular terms and call their sum $A_k=\sum_{i\le k}\sigma_iu_iv_i^*$. Terms beyond the rank are zero. For $k\le r$, storing the vectors and coefficients costs $k(m+n+1)$ numbers, rather than $mn$. What error does this saving force?

\begin{theorem}[Schmidt, Eckart--Young, Mirsky]\label{linear:thm:eckartyoung}
Let $0\le k\le p=\min(m,n)$, and read $\sigma_{p+1}$ as zero. Every matrix $B$ of rank at most $k$ satisfies
\[
\|A-B\|\ge\sigma_{k+1}=\|A-A_k\|,\qquad
\|A-B\|_F^2\ge\sum_{i>k}\sigma_i^2=\|A-A_k\|_F^2.
\]
Thus truncation gives a best approximation of rank at most $k$ in each of these two norms.
\end{theorem}
\begin{proof}
The omitted singular terms are an SVD of $A-A_k$, giving both equalities. If $k=p$, that difference is zero and both lower bounds are automatic. Assume $k<p$.

For the operator bound, $\dim\ker B\ge n-k$, whereas $\operatorname{span}(v_1,\ldots,v_{k+1})$ has dimension $k+1$. Their intersection contains a unit vector $x=\sum_{i\le k+1}c_iv_i$. Since $Bx=0$,
\[
\|(A-B)x\|^2=\|Ax\|^2=\sum_{i\le k+1}\sigma_i^2|c_i|^2\ge\sigma_{k+1}^2.
\]
For the Frobenius bound, apply~\eqref{linear:eq:weylsv} to $A=(A-B)+B$ with $j=k+1$. The rank assumption gives $\sigma_{k+1}(B)=0$, and therefore $\sigma_{i+k}(A)\le\sigma_i(A-B)$ whenever $i+k\le p$. Square and sum these inequalities:
\[
\sum_{j>k}\sigma_j(A)^2\le\sum_{i=1}^{p-k}\sigma_i(A-B)^2\le\|A-B\|_F^2.
\]
\end{proof}

\begin{remark}
Mirsky's 1960 theorem gives the same minimiser for every unitarily invariant norm. One route uses Ky Fan dominance: all such norms satisfy $\|C\|\le\|D\|$ exactly when every Ky Fan sum for $C$ is at most the corresponding sum for $D$. The preceding proof actually gives termwise domination $\sigma_i(A-A_k)\le\sigma_i(A-B)$, with zeros appended, and hence all those sum inequalities. Bhatia's book in the sources proves the dominance theorem; we do not need it for either norm proved here.
\end{remark}

\begin{figure}[ht]
\centering
\fig{linear-ch06-lowrank}
\caption{A $32\times32$ picture and its best approximations of rank $1$, $3$ and $8$. The relative errors in Frobenius norm,
$\bigl(\sum_{i>k}\sigma_i^2/\sum_i\sigma_i^2\bigr)^{1/2}$, are about $0.30$, $0.16$ and $0.07$. Rank $8$ needs $8\cdot65=520$
numbers instead of $1024$.}
\label{linear:fig:lowrank}
\end{figure}

For the opening photograph, a singular triple costs $2001$ numbers. Twenty-four cost $48{,}024$; twenty-five cost $50{,}025$ and exceed the budget. No matrix of rank at most $24$ has smaller operator or Frobenius error than $A_{24}$. These mathematical error measures do not fully describe human vision; image formats may optimise other features.

For a square invertible matrix, take $k=n-1$. The theorem says that $\sigma_n$ is exactly the distance to the nearest singular matrix in operator norm. Losing invertibility means losing the last surviving stretch.

\section{The pseudo-inverse and least squares}

An inverse would undo every stretch. If some stretches vanish, undo the positive ones and send the unused directions to zero.

\begin{definition}\label{linear:def:pinv}
For an SVD $A=\sum_{i\le r}\sigma_iu_iv_i^*$, define the \emph{Moore--Penrose pseudo-inverse} by
\[
A^+=\sum_{i\le r}\sigma_i^{-1}v_iu_i^*\in M_{n,m}(\mathbb K).
\]
\end{definition}

The next result also shows that this formula is independent of the chosen singular bases: its value on each $b$ is uniquely specified by a minimisation problem.

\begin{theorem}\label{linear:thm:pinv}
For each $b\in\mathbb K^m$, the vector $x^+=A^+b$ minimises $\|Ax-b\|$. Among all vectors attaining that minimum, it is the unique vector of smallest norm.
\end{theorem}
\begin{proof}
Expand $x=\sum_jy_jv_j$ and $b=\sum_ic_iu_i$. Orthogonality separates the error into
\[
\|Ax-b\|^2=\sum_{i\le r}|\sigma_iy_i-c_i|^2+\sum_{i>r}|c_i|^2.
\]
The second sum is fixed. The first vanishes exactly when $y_i=c_i/\sigma_i$ for $i\le r$. The other input coordinates are free. Since $\|x\|^2=\sum_j|y_j|^2$, its minimum among these solutions occurs uniquely when every free coordinate is zero. That vector is $\sum_{i\le r}\sigma_i^{-1}c_iv_i=A^+b$.
\end{proof}

If the columns of $A$ are independent, singular coordinates give $A^+=(A^*A)^{-1}A^*$. This recovers the usual normal-equation formula of Theorem~\ref{linear:thm:lsq} without assuming it in the proof. For an invertible square matrix the formula reduces to $A^{-1}$. Penrose's four equations give another intrinsic description; Exercises~\ref{linear:ex:pinv} and~\ref{linear:ex:penroseunique} prove it.

\section{Principal component analysis}

Suppose the rows of $X\in M_{N,n}(\R)$ are points $x_1,\ldots,x_N$ with $\sum_ix_i=0$. We want a $k$-dimensional subspace $W$ for which the total squared orthogonal distance $\sum_id(x_i,W)^2$ is smallest. Centring matters: without it we would also have to choose the location of an affine plane. Pearson asked about lines and planes of closest fit in 1901; Hotelling's 1933 work developed principal components.

\begin{theorem}[principal components]\label{linear:thm:pca}
Write an SVD of $X$ as $X=\sum_i\sigma_iu_iv_i^{\mathsf T}$, with a complete right singular basis. For $0\le k\le n$, the space $W_k=\operatorname{span}(v_1,\ldots,v_k)$ attains the minimum of $\sum_id(x_i,W)^2$ over all $k$-dimensional subspaces $W$. The minimum is $\sum_{j>k}\sigma_j^2$, with zero singular values appended when needed.
\end{theorem}
\begin{proof}
Let $P_W$ be orthogonal projection onto $W$. The rows of $XP_W$ are the projected data points, so the quantity to minimise is $\|X-XP_W\|_F^2$. The matrix $XP_W$ has rank at most $k$. If $k<\min(N,n)$, Theorem~\ref{linear:thm:eckartyoung} bounds this error below by $\sum_{j>k}\sigma_j^2$; if $k\ge\min(N,n)$, the same lower bound is zero. For $W_k$, multiplication by its projection keeps precisely the first $k$ right singular directions. Thus $XP_{W_k}=X_k$, and the lower bound is attained in every case.
\end{proof}

The vectors $v_i$ are the \emph{principal directions}. They diagonalise the covariance matrix $N^{-1}X^{\mathsf T}X$, and $\sigma_i^2/N$ is the variance in direction $v_i$. When the data are not all zero, the fraction $\sum_{i\le k}\sigma_i^2/\sum_i\sigma_i^2$ records how much variance the first $k$ directions retain. Least squares regression asks a different question: it singles out a dependent coordinate and minimises vertical distances. Exercise~\ref{linear:ex:pcatoy} compares the two fitted lines explicitly.

\section{Conditioning}

Take an invertible square matrix $A$. Perturb the right-hand side of $Ax=b$ and ask how much the solution can change. The relevant ratio is the \emph{condition number}
\[
\kappa(A)=\|A\|\|A^{-1}\|=\frac{\sigma_1}{\sigma_n}.
\]

\begin{proposition}\label{linear:prop:condition}
Suppose $Ax=b\ne0$ and $A(x+\delta x)=b+\delta b$. Then
\[
\frac{\|\delta x\|}{\|x\|}\le\kappa(A)\frac{\|\delta b\|}{\|b\|}.
\]
Suitable choices of $b$ and $\delta b$ attain equality. Furthermore,
\[
\min\left\{\frac{\|E\|}{\|A\|}: A+E\text{ is singular}\right\}=\frac{\sigma_n}{\sigma_1}=\frac1{\kappa(A)}.
\]
\end{proposition}
\begin{proof}
Subtract the two equations to get $\delta x=A^{-1}\delta b$. Consequently $\|\delta x\|\le\sigma_n^{-1}\|\delta b\|$, while $\|b\|\le\sigma_1\|x\|$. Multiply the two bounds. To attain equality, choose $x=v_1$, so $b=\sigma_1u_1$, and take $\delta b$ along $u_n$. Finally, Theorem~\ref{linear:thm:eckartyoung} with $k=n-1$ gives distance $\sigma_n$ to singular matrices. Divide by $\|A\|=\sigma_1$.
\end{proof}

A condition number $10^d$ permits an amplification of relative error by $10^d$. Thus errors of size $10^{-16}$ can leave only about $16-d$ decimal digits in a worst-case estimate. This is a sensitivity statement about the problem; the algorithm can introduce further errors. Orthogonal matrices have condition number $1$. Forming normal equations replaces the condition number of a full-column-rank matrix $A$ by its square, whereas its $QR$ factor $R$ keeps the original singular values. Exercise~\ref{linear:ex:normal2x2} checks both facts. The Hilbert matrices in Appendix~\ref{linear:ch:inner} give striking examples of ill conditioning.

\section{Procrustes}

Given two point configurations, can a rigid change of orientation bring them close? Let their coordinates be the columns of $A,B\in M_{n,N}(\R)$. We allow an orthogonal transformation, hence either a rotation or a reflection. The name \emph{Procrustes} recalls the mythical host who forced travellers to fit his bed. Here the lengths within the first configuration stay intact.

\begin{theorem}[orthogonal Procrustes]\label{linear:thm:procrustes}
Choose an SVD $BA^{\mathsf T}=U\Sigma V^{\mathsf T}$. Among all orthogonal matrices $Q$, the choice $Q=UV^{\mathsf T}$ minimises $\|QA-B\|_F$.
\end{theorem}
\begin{proof}
Expand the square:
\[
\|QA-B\|_F^2=\|A\|_F^2+\|B\|_F^2-2\operatorname{tr}(B^{\mathsf T}QA).
\]
Only the trace depends on $Q$. Cyclically permuting its factors gives
\[
\operatorname{tr}(B^{\mathsf T}QA)=\operatorname{tr}(QAB^{\mathsf T})
=\operatorname{tr}(U^{\mathsf T}QV\Sigma).
\]
Put $Z=U^{\mathsf T}QV$. It is orthogonal, so every diagonal entry satisfies $z_{ii}\le1$. Therefore the trace is $\sum_i\sigma_iz_{ii}\le\sum_i\sigma_i$. Taking $Z=I$, or $Q=UV^{\mathsf T}$, attains the upper bound and hence the smallest squared distance.
\end{proof}

For a square real $A$, this also finds a nearest orthogonal matrix to $A$ itself: it is a unitary factor in its polar decomposition. Exercise~\ref{linear:ex:nearestorth} checks the transpose needed to turn the configuration problem into that distance problem. Singular matrices can have more than one nearest orthogonal matrix, just as their polar unitary factor need not be unique.

\begin{beauty}
A singular vector pair records an input direction, an output direction and the stretch between them. The same pair has several concrete meanings. In a picture it is one layer of an approximation. In a data table it is one direction of variation. In a linear system the smallest stretch marks the direction in which errors can grow most. These interpretations all come from the single identity $Av_i=\sigma_iu_i$.
\end{beauty}

\section{From stretches to graphs}

The photograph's storage budget led to a theorem about every rectangular matrix. Keeping its first $24$ singular terms is optimal within that rank, but other questions use the decomposition differently: invert the positive stretches for least squares, or use the two singular bases to align point configurations. The bases make each optimisation explicit.

A graph gives us a square matrix with a more restrictive origin: its entries record which vertices are joined. How much of that discrete picture can its eigenvalues recover? The next chapter begins with colourings and independent sets, and reaches Huang's sensitivity argument.

\begin{figure}[ht]
\centering\fig{ch04-rw-budget}
\caption{For singular values $8,4,2,1$, a rank-$k$ truncation leaves the squared Frobenius error shown in coral. Each retained direction removes its own squared singular value from the error.}
\label{fig:v8c04budget}
\end{figure}
\begin{figure}[ht]
\centering\fig{ch04-rw-conditioning}
\caption{The map $A=\operatorname{diag}(4,1)$ sends equal input lengths to stretches $4$ and $1$. With $x=e_1$ and $\delta b=\varepsilon e_2$, the relative solution error is four times the relative data error.}
\label{fig:v8c04conditioning}
\end{figure}

\section{Exercises}

\begin{exercise}[$\star$]\label{linear:ex:svd2}
Find a singular value decomposition and a polar decomposition of $A=\begin{pmatrix}3&0\\4&5\end{pmatrix}$, the matrix in Figure~\ref{linear:fig:ellipse}. Verify both identities $\sigma_1\sigma_2=|\det A|$ and $\sigma_1^2+\sigma_2^2=\|A\|_F^2$.
\end{exercise}
\begin{exercise}[$\star$]\label{linear:ex:normsvd}
Prove that a square matrix satisfies $|\det A|=\prod_i\sigma_i$. Bound the modulus of each eigenvalue by $\sigma_1$. Give a $2\times2$ example for which no eigenvalue equals a singular value. For a normal matrix, prove that the singular values are the absolute values of the eigenvalues.
\end{exercise}
\begin{exercise}[$\star$]\label{linear:ex:normineq}
With $r=\operatorname{rank}A$, establish $\|A\|\le\|A\|_F\le\sqrt r\,\|A\|$, including the case $r=0$. For compatible matrices also establish $\|AB\|_F\le\|A\|\|B\|_F$.
\end{exercise}
\begin{exercise}[$\star$]\label{linear:ex:pinv}
Verify the four Penrose equations for $X=A^+$:
\[
AXA=A,\quad XAX=X,\quad (AX)^*=AX,\quad (XA)^*=XA.
\]
Identify $AA^+$ and $A^+A$ as orthogonal projections onto $\operatorname{Im}A$ and $\operatorname{Im}A^*$, respectively. Find $A^+$ when every entry of the $2\times2$ matrix $A$ is $1$.
\end{exercise}
\begin{exercise}[$\star\star$]\label{linear:ex:penroseunique}
Prove uniqueness in Penrose's characterisation: a matrix $X$ satisfying all four equations in Exercise~\ref{linear:ex:pinv} must equal $A^+$.
\end{exercise}
\begin{exercise}[$\star\star$]\label{linear:ex:pcatoy}
For the points $(-2,-1),(-1,-1),(0,0),(1,1),(2,1)$, calculate the first principal direction. Compare its slope with the regression slope in $y=bx$. Express the principal slope using $\varphi=(1+\sqrt5)/2$.
\end{exercise}
\begin{exercise}[$\star\star$]\label{linear:ex:nearestorth}
Let $A$ be square and real, and choose its polar decomposition $A=WP$. Prove that $W$ is a nearest orthogonal matrix in Frobenius norm and that its squared distance is $\sum_i(\sigma_i-1)^2$.
\end{exercise}
\begin{exercise}[$\star\star$]\label{linear:ex:nuclear}
For square matrices over $\mathbb K$, prove
$\|A\|_*=\max_{Q\text{ unitary}}|\operatorname{tr}(AQ)|$.
Use this to obtain $|\operatorname{tr}(AB)|\le\|A\|_*\|B\|$ for every square $B$ of the same size.
\end{exercise}
\begin{exercise}[$\star\star$]\label{linear:ex:svinterlace}
Delete one column of $A\in M_{m,n}(\mathbb K)$ to obtain $B$. Show that their singular values interlace:
$\sigma_1(A)\ge\sigma_1(B)\ge\sigma_2(A)\ge\sigma_2(B)\ge\cdots$, for all indices present, with zeros appended where necessary.
\end{exercise}
\begin{exercise}[$\star\star$]\label{linear:ex:svprod}
For compatible $A,B$, prove $\sigma_i(AB)\le\|A\|\sigma_i(B)$ at every available index, using zero extensions if needed. Conclude that square invertible matrices obey $\kappa(AB)\le\kappa(A)\kappa(B)$.
\end{exercise}
\begin{exercise}[$\star\star$]\label{linear:ex:normal2x2}
Assume $A$ has independent columns and define $\kappa(A)=\sigma_1/\sigma_n$. Verify $\kappa(A^*A)=\kappa(A)^2$. If $A=QR$ with $Q^*Q=I$, verify instead $\kappa(R)=\kappa(A)$.
\end{exercise}
\begin{exercise}[$\star\star\star$]\label{linear:ex:matrixcompletion}
A rank-one $3\times3$ matrix has known entries $a_{11}=2$, $a_{12}=4$, $a_{21}=3$, $a_{33}=5$, $a_{23}=10$. Determine every rank-one completion. Then consider an $m\times n$ rank-one matrix with no zero entries. Prove that observing a whole row and a whole column determines it, and that fewer than $m+n-1$ observations cannot ensure uniqueness in general.
\end{exercise}

\begin{exercise}[$\star$]\label{ex:v8c04a}
For $A=\begin{pmatrix}2&0&0\\0&1&0\end{pmatrix}$, find $A^+$ and both associated orthogonal projections. For $b=(4,3)$, describe all least squares solutions and identify the one of smallest norm.
\end{exercise}
\begin{exercise}[$\star$]\label{ex:v8c04b}
A matrix has singular values $8,4,2,1$. Find the smallest rank whose best approximation has relative Frobenius error below $1/4$. For a $100\times80$ matrix with these singular values and all others zero, compare the storage of that approximation with all $8000$ entries, counting singular vectors and values separately as in the text.
\end{exercise}
\begin{exercise}[$\star\star$]\label{ex:v8c04c}
Let $A=\operatorname{diag}(4,1)$, $b=(4,0)$ and $\delta b=(0,\varepsilon)$ with $\varepsilon\ne0$. Compute both relative errors and show that the condition bound is attained. Find a perturbation $E$ of minimum operator norm making $A+E$ singular.
\end{exercise}
\begin{exercise}[$\star\star$]\label{ex:v8c04d}
For centred points $(2,0),(-2,0),(0,1),(0,-1)$, determine the best fitting line through zero, its total squared orthogonal error and its explained variance fraction. Find every line attaining that minimum.
\end{exercise}

\section{Hints}

\begin{hint}{linear:ex:svd2}
Begin by diagonalising the Gram matrix $\begin{pmatrix}25&20\\20&25\end{pmatrix}$.
\end{hint}
\begin{hint}{linear:ex:normsvd}
Take determinants in $A^*A$. Apply $A$ to a unit eigenvector for the norm bound. An upper triangular shear gives the requested example.
\end{hint}
\begin{hint}{linear:ex:normineq}
Only $r$ singular values are non-zero. For the product estimate, sum the squared norms of the images of the columns of $B$.
\end{hint}
\begin{hint}{linear:ex:pinv}
The products become diagonal projections in singular coordinates. For the example use the unit vector $u=(1,1)/\sqrt2$ and $A=2uu^{\mathsf T}$.
\end{hint}
\begin{hint}{linear:ex:penroseunique}
Compare two solutions $X,Y$. First show that the Hermitian idempotents $AX,AY$ are the same orthogonal projection; then compare $XA,YA$.
\end{hint}
\begin{hint}{linear:ex:pcatoy}
The two fitted lines use the same sums, $\sum x^2=10$, $\sum xy=6$ and $\sum y^2=4$, in different ways.
\end{hint}
\begin{hint}{linear:ex:nearestorth}
Apply Procrustes to $A$ and the target $I$, remembering that $\|QA-I\|_F=\|A-Q^{\mathsf T}\|_F$.
\end{hint}
\begin{hint}{linear:ex:nuclear}
For unitary $Q$, cycle the trace into singular coordinates. For arbitrary $B$, those coordinates still bound every diagonal entry by $\|B\|$.
\end{hint}
\begin{hint}{linear:ex:svinterlace}
Compare the eigenvalues of the two Gram matrices; one is a principal submatrix of the other.
\end{hint}
\begin{hint}{linear:ex:svprod}
Compare $\|ABx\|^2$ and $\|A\|^2\|Bx\|^2$ before applying the min--max formula.
\end{hint}
\begin{hint}{linear:ex:normal2x2}
The Gram matrix has eigenvalues $\sigma_i^2$. For $QR$, the identity $Q^*Q=I$ leaves the Gram matrix unchanged.
\end{hint}
\begin{hint}{linear:ex:matrixcompletion}
Write each entry as $x_iy_j$. Represent an observation by an edge between its row vertex and its column vertex.
\end{hint}
\begin{hint}{ex:v8c04a}
The third input coordinate has no effect on the output. Invert the two positive diagonal entries and set that coordinate to zero.
\end{hint}
\begin{hint}{ex:v8c04b}
Compare the tail sum of squares with $85/16$; the total squared Frobenius norm is $85$.
\end{hint}
\begin{hint}{ex:v8c04c}
The original right-hand side lies in the direction of greatest stretch, while its perturbation lies in the direction of least stretch.
\end{hint}
\begin{hint}{ex:v8c04d}
For a unit line direction $w=(\cos\theta,\sin\theta)$, the squared length captured from all four points is $8\cos^2\theta+2\sin^2\theta$.
\end{hint}

\section{Solutions}

\begin{solution}{linear:ex:svd2}
\idea Use the eigenvectors of $A^{\mathsf T}A$ as input directions, then normalise their images.
\writeup The Gram matrix has eigenvalues $45,5$ and unit eigenvectors $v_1=(1,1)/\sqrt2$, $v_2=(-1,1)/\sqrt2$. Hence $\sigma_1=3\sqrt5$, $\sigma_2=\sqrt5$, and
\[
u_1=\frac{Av_1}{3\sqrt5}=\frac{(1,3)}{\sqrt{10}},\qquad
u_2=\frac{Av_2}{\sqrt5}=\frac{(-3,1)}{\sqrt{10}}.
\]
Taking these columns for $U,V$ gives the SVD. The requested checks read $\sigma_1\sigma_2=15=|\det A|$ and $\sigma_1^2+\sigma_2^2=50=\|A\|_F^2$.
For the polar factors, direct multiplication gives
\[
P=V\begin{pmatrix}3\sqrt5&0\\0&\sqrt5\end{pmatrix}V^{\mathsf T}
=\sqrt5\begin{pmatrix}2&1\\1&2\end{pmatrix},\qquad
W=UV^{\mathsf T}=\frac1{\sqrt5}\begin{pmatrix}2&-1\\1&2\end{pmatrix}.
\]
The factor $W$ rotates by $\arctan(1/2)$, and $WP=\begin{pmatrix}3&0\\4&5\end{pmatrix}$ confirms the factorisation.
\end{solution}

\begin{solution}{linear:ex:normsvd}
\idea The Gram eigenvalues give the determinant, while a single eigenvector gives the eigenvalue bound.
\writeup Taking determinants yields $|\det A|^2=\det(A^*A)=\prod_i\sigma_i^2$. Non-negativity allows square roots. If $Ax=\lambda x$ and $\|x\|=1$, then $|\lambda|=\|Ax\|\le\|A\|=\sigma_1$.
For $A=\begin{pmatrix}1&1\\0&1\end{pmatrix}$ the eigenvalues are both $1$, but the Gram eigenvalues are $(3\pm\sqrt5)/2$. Its singular values are $(\sqrt5+1)/2$ and $(\sqrt5-1)/2$, neither equal to $1$.
Finally, normality gives $A=WDW^*$. Write $D=\Theta|D|$ with $\Theta$ diagonal unitary, choosing phase $1$ at zero entries. Then $A=(W\Theta)|D|W^*$ is an SVD after reordering, which proves the asserted moduli formula.
\end{solution}

\begin{solution}{linear:ex:normineq}
\idea Compare the largest squared stretch with the sum of all squared stretches.
\writeup If $r=0$, the matrix is zero and both inequalities are equalities. Otherwise
\[
\sigma_1^2\le\sum_{i=1}^r\sigma_i^2\le r\sigma_1^2.
\]
Taking square roots proves the two norm comparisons. Write $b_j$ for the columns of $B$. The product has columns $Ab_j$, so
$\|AB\|_F^2=\sum_j\|Ab_j\|^2\le\|A\|^2\sum_j\|b_j\|^2=\|A\|^2\|B\|_F^2$.
\end{solution}

\begin{solution}{linear:ex:pinv}
\idea Multiplying the two singular expansions leaves one orthogonal projection on each side.
\writeup Orthogonality gives
\[
AA^+=\sum_{i\le r}u_iu_i^*,\qquad A^+A=\sum_{i\le r}v_iv_i^*.
\]
These are the projections onto the image spaces specified in the question, so they are Hermitian. The first projection fixes every column of $A$, giving $AA^+A=A$; the second fixes every column of $A^+$, giving $A^+AA^+=A^+$. These are the other two equations.
For the matrix of ones, $A=2uu^{\mathsf T}$ with $u=(1,1)/\sqrt2$. Its single positive singular value is $2$, whence $A^+=\tfrac12uu^{\mathsf T}=\tfrac14\begin{pmatrix}1&1\\1&1\end{pmatrix}$.
\end{solution}

\begin{solution}{linear:ex:penroseunique}
\idea Orthogonal projections are determined by their image, or by their kernel.
\writeup Suppose $X$ satisfies the four equations. Then $AX$ is Hermitian and $(AX)^2=AXAX=AX$. Its image is contained in $\operatorname{Im}A$, and $AX$ fixes that whole image because $AXA=A$. Thus $AX$ is its orthogonal projection. Also $XA$ is Hermitian and idempotent. Its kernel equals $\ker A$: inclusion from right to left is immediate, and $XAz=0$ implies $Az=AXAz=0$. Hence $XA$ is projection onto $(\ker A)^\perp$.
If $Y$ is another solution, these descriptions give $AX=AY$ and $XA=YA$. Therefore
\[
X=XAX=X(AY)=(XA)Y=(YA)Y=YAY=Y.
\]
Since $A^+$ satisfies the equations, it is the only solution.
\end{solution}

\begin{solution}{linear:ex:pcatoy}
\idea Find a leading Gram eigenvector for PCA and minimise a scalar quadratic for regression.
\writeup Both coordinate sums vanish, so the data are centred. Their Gram matrix is $G=\begin{pmatrix}10&6\\6&4\end{pmatrix}$, whose trace is $14$ and determinant $4$. Its eigenvalues are $7\pm3\sqrt5$. The vector $(6,3\sqrt5-3)$ belongs to the larger one; its slope is $(\sqrt5-1)/2=1/\varphi$.
The regression error is $\sum_i(y_i-bx_i)^2=4-12b+10b^2$, minimised at $b=3/5$. Thus the orthogonal fit has slope $1/\varphi$, while the vertical fit has slope $3/5$. The first principal component retains the variance fraction $(7+3\sqrt5)/14$.
\end{solution}

\begin{solution}{linear:ex:nearestorth}
\idea Transpose the fitted orthogonal matrix after applying Procrustes with target $I$.
\writeup If $A=U\Sigma V^{\mathsf T}$, then $A^{\mathsf T}=V\Sigma U^{\mathsf T}$. Procrustes chooses $Q=VU^{\mathsf T}$ to minimise $\|QA-I\|_F$. Orthogonal invariance gives $\|QA-I\|_F=\|A-Q^{\mathsf T}\|_F$, so $W=UV^{\mathsf T}$ is nearest to $A$. Moreover $A-W=U(\Sigma-I)V^{\mathsf T}$, and its squared Frobenius norm is exactly $\sum_i(\sigma_i-1)^2$. Every polar unitary factor arises from an SVD with a suitable choice on the zero singular subspaces, so the conclusion applies to the chosen $W$ even when $A$ is singular.
\end{solution}

\begin{solution}{linear:ex:nuclear}
\idea In singular coordinates, bound each diagonal entry of the second factor by its operator norm.
\writeup For $A=U\Sigma V^*$ and unitary $Q$, cyclicity of trace gives
\[
|\operatorname{tr}(AQ)|=|\operatorname{tr}(\Sigma V^*QU)|
\le\sum_i\sigma_i |(V^*QU)_{ii}|\le\sum_i\sigma_i.
\]
The choice $Q=VU^*$ attains equality. For arbitrary $B$, put $C=V^*BU$. Every entry $|c_{ii}|=|\langle Ce_i,e_i\rangle|$ is at most $\|C\|=\|B\|$ by Cauchy--Schwarz. The same trace calculation now gives $|\operatorname{tr}(AB)|\le\sum_i\sigma_i\|B\|$. This proof works over both fields without introducing complex unitary matrices into a real argument.
\end{solution}

\begin{solution}{linear:ex:svinterlace}
\idea Deleting an input coordinate gives a principal compression of the Gram matrix.
\writeup If column $j$ is removed, $B^*B$ is obtained from $A^*A$ by removing row and column $j$. Theorem~\ref{linear:thm:interlacing} gives
$\lambda_k(A^*A)\ge\lambda_k(B^*B)\ge\lambda_{k+1}(A^*A)$.
All these numbers are non-negative. Their square roots are the singular values, padded by the zero Gram eigenvalues when a matrix is wider than it is tall. Taking square roots proves the interlacing at every indicated index.
\end{solution}

\begin{solution}{linear:ex:svprod}
\idea A pointwise quadratic-form bound passes through the min--max characterisation.
\writeup For every input vector $x$,
\[
x^*(AB)^*ABx=\|ABx\|^2\le\|A\|^2\|Bx\|^2=\|A\|^2x^*B^*Bx.
\]
Apply Theorem~\ref{linear:thm:courantfischer} to these two Hermitian matrices. It gives the same inequality for their ordered eigenvalues; square roots give $\sigma_i(AB)\le\|A\|\sigma_i(B)$. If an index exceeds the rank bounds of either factor, the relevant singular values vanish as required. For invertible square factors, submultiplicativity gives
\[
\|AB\|\le\|A\|\|B\|,\qquad
\|(AB)^{-1}\|=\|B^{-1}A^{-1}\|\le\|B^{-1}\|\|A^{-1}\|.
\]
Multiplying proves the condition-number bound. Applying the first part to $B^*A^*$ also gives $\sigma_i(AB)\le\sigma_i(A)\|B\|$, with the same zero-extension convention.
\end{solution}

\begin{solution}{linear:ex:normal2x2}
\idea Both assertions can be read from one Gram matrix.
\writeup Independent columns give $n$ positive singular values. The matrix $A^*A=V(\Sigma^*\Sigma)V^*$ is positive definite with eigenvalues, and hence singular values, $\sigma_1^2,\ldots,\sigma_n^2$. Its condition number is therefore $\sigma_1^2/\sigma_n^2$.
For $A=QR$ with $Q^*Q=I$, compute $A^*A=R^*Q^*QR=R^*R$. The identical Gram spectra give identical singular values and $\kappa(R)=\kappa(A)$. This comparison explains the sensitivity cost of normal equations; it does not by itself certify every implementation of a solver.
\end{solution}

\begin{solution}{linear:ex:matrixcompletion}
\idea Products $x_iy_j$ propagate across observed entries; a disconnected observation pattern leaves a scaling freedom.
\writeup Write $A=xy^{\mathsf T}$. The first row observations give $y_2=2y_1$, and the first column gives $x_2=3x_1/2$, with $x_1y_1=2$. Since $x_2y_3=10$, we have $y_3=20/(3x_1)$. The last observation gives $x_3=3x_1/4$. Substitution determines every remaining entry:
\[
A=\begin{pmatrix}2&4&20/3\\3&6&10\\3/2&3&5\end{pmatrix}
=\begin{pmatrix}1\\3/2\\3/4\end{pmatrix}\begin{pmatrix}2&4&20/3\end{pmatrix}.
\]
The displayed factorisation proves existence as well as uniqueness.
In general, non-zero row $i$ and column $j$ determine all entries by $a_{kl}=a_{kj}a_{il}/a_{ij}$. For the lower bound on observations, make a bipartite graph with one vertex for each row and each column. Join them at observed positions. If this graph is disconnected, choose a connected component containing a row and having some column outside it, or a component containing a column and having some row outside it. Such a choice exists because both types of vertices occur. Multiply its row factors by $t\ne0$ and its column factors by $t^{-1}$. Every observed product stays fixed, while a product crossing to the indicated outside vertex changes when $t\ne1$. Hence uniqueness requires connectedness. A connected graph on $m+n$ vertices has at least $m+n-1$ edges, giving the bound. For higher-rank completion under random sampling, see Cand\`es and Recht in the sources; the rank-one argument alone does not prove that theory.
\end{solution}

\begin{solution}{ex:v8c04a}
\idea Solve the two active coordinates and leave the kernel coordinate free.
\writeup Directly,
\[
A^+=\begin{pmatrix}1/2&0\\0&1\\0&0\end{pmatrix},\quad
AA^+=I_2,\quad A^+A=\operatorname{diag}(1,1,0).
\]
The equation $Ax=b$ reads $2x_1=4$, $x_2=3$. Thus all minimisers, with zero residual, are $(2,3,t)$ for $t\in\mathbb K$. Their squared norms are $13+|t|^2$, so the unique smallest one is $A^+b=(2,3,0)$, of norm $\sqrt{13}$.
\end{solution}

\begin{solution}{ex:v8c04b}
\idea Relative Frobenius error is the square root of the discarded energy fraction.
\writeup The total squared norm is $64+16+4+1=85$. At rank one the discarded sum is $21$, giving relative error $\sqrt{21/85}>1/4$ because $21>85/16$. At rank two it is $5$, and $5<85/16$, so $\sqrt{5/85}=1/\sqrt{17}<1/4$. The smallest rank is therefore two. Each singular triple for a $100\times80$ matrix stores $181$ numbers. Two triples use $362$ numbers instead of $8000$; the orthogonality constraints could allow other encodings, but are not part of this count.
\end{solution}

\begin{solution}{ex:v8c04c}
\idea Choose the data and its perturbation along opposite extreme singular directions.
\writeup We have $x=(1,0)$ and $\delta x=(0,\varepsilon)$. Hence $\|\delta x\|/\|x\|=|\varepsilon|$, while $\|\delta b\|/\|b\|=|\varepsilon|/4$. Their ratio is $4=\kappa(A)$. Taking $E=\operatorname{diag}(0,-1)$ makes $A+E=\operatorname{diag}(4,0)$ singular and has norm $1$. No smaller perturbation can do so, since the distance theorem gives minimum norm $\sigma_2(A)=1$.
\end{solution}

\begin{solution}{ex:v8c04d}
\idea Maximise the total squared projection onto the line.
\writeup The Gram matrix of the four centred rows is $\operatorname{diag}(8,2)$ and the total squared distance to zero is $10$. A unit direction $w=(\cos\theta,\sin\theta)$ captures $8\cos^2\theta+2\sin^2\theta=2+6\cos^2\theta$. Its maximum is $8$, attained exactly when $\sin\theta=0$. Thus the only minimising line is the $x$-axis, its residual sum of squares is $10-8=2$, and it explains the fraction $8/10=4/5$ of the variance. The two unit vectors $e_1,-e_1$ describe that same line.
\end{solution}

\section*{Sources and further reading}

\begin{itemize}
\item E.~Beltrami, \emph{Sulle funzioni bilineari}, Giornale di Matematiche ad uso degli studenti delle universit\`a 11
(1873), 98--106; C.~Jordan, \emph{M\'emoire sur les formes bilin\'eaires}, Journal de math\'ematiques pures et appliqu\'ees
19 (1874), 35--54.
\item G.~W.~Stewart, \emph{On the early history of the singular value decomposition}, SIAM Review 35 (1993), 551--566.
\item E.~Schmidt, \emph{Zur Theorie der linearen und nichtlinearen Integralgleichungen I}, Mathematische Annalen 63 (1907),
433--476; C.~Eckart and G.~Young, \emph{The approximation of one matrix by another of lower rank}, Psychometrika 1 (1936),
211--218; L.~Mirsky, \emph{Symmetric gauge functions and unitarily invariant norms}, Quarterly Journal of Mathematics 11
(1960), 50--59.
\item K.~Pearson, \emph{On lines and planes of closest fit to systems of points in space}, Philosophical Magazine 2 (1901),
559--572; H.~Hotelling, \emph{Analysis of a complex of statistical variables into principal components}, Journal of
Educational Psychology 24 (1933), 417--441.
\item R.~Penrose, \emph{A generalized inverse for matrices}, Proceedings of the Cambridge Philosophical Society 51 (1955),
406--413.
\item G.~Golub and W.~Kahan, \emph{Calculating the singular values and pseudo-inverse of a matrix}, SIAM Journal on
Numerical Analysis 2 (1965), 205--224.
\item P.~H.~Sch\"onemann, \emph{A generalized solution of the orthogonal Procrustes problem}, Psychometrika 31 (1966), 1--10.
\item Y.~Koren, R.~Bell and C.~Volinsky, \emph{Matrix factorization techniques for recommender systems}, IEEE Computer 42
(2009), 30--37; E.~J.~Cand\`es and B.~Recht, \emph{Exact matrix completion via convex optimization}, Foundations of
Computational Mathematics 9 (2009), 717--772.
\item R.~Bhatia, \emph{Matrix Analysis}, Springer, 1997, Chapter~IV (unitarily invariant norms, Ky Fan dominance); L.~N.~Trefethen
and D.~Bau, \emph{Numerical Linear Algebra}, SIAM, 1997, Lectures 4--5.
\end{itemize}
